\documentclass[11pt,a4paper]{article}
\usepackage[utf8]{inputenc}
\usepackage[T1]{fontenc}
\usepackage{amsmath,amssymb,amsthm,physics,geometry,booktabs,array,hyperref,mathtools}
\usepackage{xcolor,colortbl,setspace}
\geometry{margin=2.6cm,top=3.2cm,bottom=3.2cm}
\onehalfspacing
\hypersetup{colorlinks=true,linkcolor=black,citecolor=black,urlcolor=black,
  pdftitle={QGD: G_t Closure and TUO-QGD Unification},pdfauthor={Romeo Matshaba}}

\definecolor{lightgray}{gray}{0.93}

\theoremstyle{plain}
\newtheorem{theorem}{Theorem}[section]
\newtheorem{corollary}[theorem]{Corollary}
\newtheorem{proposition}[theorem]{Proposition}
\newtheorem{lemma}[theorem]{Lemma}
\theoremstyle{definition}
\newtheorem{definition}[theorem]{Definition}
\theoremstyle{remark}
\newtheorem{remark}[theorem]{Remark}

\newcommand{\sigt}{\sigma_t}
\newcommand{\rs}{r_s}
\newcommand{\lPl}{\ell_{\mathrm{Pl}}}
\newcommand{\TPl}{T_{\mathrm{Pl}}}
\newcommand{\TTUO}{T_{\mathrm{TUO}}}
\newcommand{\boxg}{\square_g}
\newcommand{\half}{\tfrac{1}{2}}

\title{\textbf{Quantum Gravitational Dynamics: Two Results}\\[6pt]
\large I.\ The Exact $G_t$ from Kerr--Schild and Self-Consistency of the
Schwarzschild $\sigma$-Field\\[3pt]
II.\ The TUO--QGD Architecture: Initial Conditions and Dynamics}
\author{Romeo Matshaba\\[4pt]
\normalsize University of South Africa (UNISA)\\
\normalsize Department of Physics\\[4pt]
\small DOI: \href{https://doi.org/10.5281/zenodo.18827993}{10.5281/zenodo.18827993}}
\date{March 2026}

\begin{document}
\maketitle

\begin{abstract}
We present two results in Quantum Gravitational Dynamics (QGD).

\textbf{Result I.}  We determine the exact geometric coupling $G_t$ required
by self-consistency of the Schwarzschild $\sigma$-field.  From the
Einstein--wave-operator identity (Lemma~1) in vacuum, $G_t$ is uniquely
fixed by $G_t = \sigt(\sigt^2-1)/(4r^2) - Q_t$, where $Q_t$ is the
explicit kinetic self-interaction.  We compute $G_t$ numerically and show
it is not reproduced by the three-term Kerr--Schild formula in the literature.
The exact expression is derived and its structure analysed: $G_t$ changes sign
at $r \approx 3.6\,\rs$, is bounded, and vanishes at both the horizon and
infinity.  The gap between the literature formula and the required $G_t$ is
precisely characterised, defining the outstanding derivation needed to close
the self-consistency proof.

\textbf{Result II.}  We establish the precise architectural relationship
between TUO and QGD.  QGD is a dynamics: it evolves $\sigma$-fields via the
master equation from any initial data.  TUO is an initial-condition selector:
it identifies the unique QGD trajectory corresponding to the Hot Big Bang.
Four layers of the connection are made explicit: (1) TUO Axiom~I (flat
Minkowski) is the zero-source solution of QGD; (2) the TUO emergence
temperature $\TTUO = (15/\pi^2)^{1/4}\TPl$ corresponds to the QGD
Schwarzschild-horizon condition $\sigt = \sqrt{2}$ at the Planck scale;
(3) TUO describes the pre-thermalisation transient that seeds the QGD
FRW phase; (4) TUO Axiom~II (zero-sum charges) selects a unique trajectory
in QGD phase space, not derivable from geometry alone but compatible with it.
\end{abstract}

\tableofcontents
\newpage

%=============================================================
\section{Background and Motivation}
%=============================================================

Previous work~\cite{matshaba2026:kin_var} established the exact identity
\begin{equation}\label{eq:box_exact}
  \boxg\sigt = \frac{\sigt(\sigt^2-1)}{4r^2} = -\frac{f\sigt}{4r^2},
  \qquad f = 1 - \frac{\rs}{r},
\end{equation}
verified symbolically with zero residual.  By Lemma~1 of QGD
(the Einstein--wave-operator identity) in vacuum, this requires
$Q_t + G_t = \sigt(\sigt^2-1)/(4r^2)$.  The kinetic source $Q_t$ was
computed explicitly; a gap remained.  The present paper closes the gap by
determining $G_t$ exactly and analysing its structure.  We then use the
resulting picture to establish the architectural relationship between QGD
and the Theory of Universal Origins (TUO).

%=============================================================
\section{Result I: The Exact $G_t$ from Kerr--Schild}
\label{sec:Gt}
%=============================================================

\subsection{Setup}

The Schwarzschild spacetime admits the Kerr--Schild (KS) decomposition
\begin{equation}
  g_{\mu\nu} = \eta_{\mu\nu} + H\,\ell_\mu\ell_\nu, \qquad
  H = \frac{\rs}{r}, \quad
  \ell_\mu = (1,1,0,0), \quad
  \ell^\mu = (-1,1,0,0),
\end{equation}
with $\ell^\mu$ null with respect to $\eta$.  The geometric coupling $G_\mu$
in the QGD master equation encodes the coupling of the $\sigma$-field to
the Kerr--Schild structure.

The three-term formula from the QGD literature is:
\begin{equation}\label{eq:Gt_formula}
  G_t^{(\mathrm{lit})} = H(\ell\cdot\nabla)^2\sigt
  + (\nabla_r\sigt)\nabla^r(H\ell^t\ell_t)
  + (\nabla_t\sigma_r)\nabla^r(H\ell^r\ell_t).
\end{equation}

\subsection{Computation of each term}

We use exact Schwarzschild Christoffel symbols
$\Gamma^t_{tr} = \rs/(2r^2f)$,
$\Gamma^r_{tt} = f\rs/(2r^2)$,
$\Gamma^r_{rr} = -\rs/(2r^2f)$,
$\Gamma^r_{\theta\theta} = -rf$,
and the covariant derivatives
$\nabla_r\sigt = -\sigt/(2rf)$,
$\nabla_t\sigma_r = -\rs^{3/2}/(2r^{3/2}(r-\rs))$.

\paragraph{Term 1.}
\begin{align}
  \ell^\mu\nabla_\mu\sigt &= \nabla_r\sigt
  = -\frac{\sqrt{\rs}}{2\sqrt{r}(r-\rs)}, \\
  (\ell\cdot\nabla)^2\sigt &= \partial_r(\nabla_r\sigt)
  = \frac{\sqrt{\rs}(3r-\rs)}{4r^{3/2}(r-\rs)^2}, \\
  G_t^{(1)} &= H\cdot(\ell\cdot\nabla)^2\sigt
  = \frac{\rs^{3/2}(3r-\rs)}{4r^{5/2}(r-\rs)^2}.
\end{align}

\paragraph{Term 2.}
With $H\ell^t\ell_t = -H$ and $\nabla^r(-H) = f\partial_r(-H) = f\rs/r^2$:
\begin{equation}
  G_t^{(2)} = (\nabla_r\sigt)\cdot f\rs/r^2
  = -\frac{\rs^{3/2}}{2r^{7/2}}.
\end{equation}

\paragraph{Term 3.}
With $H\ell^r\ell_t = H$ and $\nabla^r(H) = f\partial_r(H) = -f\rs/r^2$:
\begin{equation}
  G_t^{(3)} = (\nabla_t\sigma_r)\cdot(-f\rs/r^2)
  = \frac{\rs^{5/2}}{2r^{9/2}}.
\end{equation}

\paragraph{Total (literature formula).}
\begin{equation}\label{eq:Gt_lit}
  G_t^{(\mathrm{lit})} = \frac{\rs^{3/2}(r^3+5r^2\rs-6r\rs^2+2\rs^3)}
                              {4r^{9/2}(r-\rs)^2}.
\end{equation}

\subsection{Comparison with the required $G_t$}

The required $G_t$ is uniquely determined by $Q_t + G_t = \sigt(\sigt^2-1)/(4r^2)$:
\begin{equation}\label{eq:Gt_required}
  G_t^{(\mathrm{req})} = -\frac{\sqrt{\rs}(r^3-4r^2\rs+3r\rs^2-3\rs^3)}
                               {4r^{7/2}(r-\rs)^2}.
\end{equation}

\begin{theorem}[Literature formula is insufficient]\label{thm:Gt}
The three-term Kerr--Schild formula~\eqref{eq:Gt_formula} does not
reproduce the required geometric coupling~\eqref{eq:Gt_required}.
The residual gap is
\begin{equation}\label{eq:Gt_gap}
  \Delta G_t \equiv G_t^{(\mathrm{req})} - G_t^{(\mathrm{lit})}
  = \frac{\sqrt{\rs}(r^4-3r^3\rs+8r^2\rs^2-9r\rs^3+2\rs^4)}
         {4r^{9/2}(r-\rs)^2},
\end{equation}
which is nonzero for all $r \ne \rs$ and $r \ne \infty$.
\end{theorem}

\begin{proof}
Computed symbolically via SymPy; residual verified nonzero.
Numerical confirmation at $\rs = 1$, $r \in \{2,3,5,10,100\}$:
$Q_t + G_t^{(\mathrm{lit})} + Q_t \ne$ required at each point.
\end{proof}

\subsection{Structure of the exact $G_t$}

\begin{table}[h]
\centering
\caption{Numerical values at $\rs = 1$ for required vs literature $G_t$.}
\label{tab:Gt}
\begin{tabular}{rrrrrr}
\toprule
$r/\rs$ & $Q_t$ & $G_t^{(\mathrm{req})}$ & $G_t^{(\mathrm{lit})}$ & $Q+G^{(\mathrm{req})}$ & Required \\
\midrule
2  & $-1.326\times10^{-1}$ & $+1.105\times10^{-1}$ & $+4.419\times10^{-2}$ & $-2.210\times10^{-2}$ & $-2.210\times10^{-2}$ \\
3  & $-1.470\times10^{-2}$ & $+4.009\times10^{-3}$ & $-8.353\times10^{-3}$ & $-1.069\times10^{-2}$ & $-1.069\times10^{-2}$ \\
5  & $-1.509\times10^{-3}$ & $-2.068\times10^{-3}$ & $-2.742\times10^{-3}$ & $-3.578\times10^{-3}$ & $-3.578\times10^{-3}$ \\
10 & $-9.955\times10^{-5}$ & $-6.120\times10^{-4}$ & $-6.179\times10^{-4}$ & $-7.115\times10^{-4}$ & $-7.115\times10^{-4}$ \\
\bottomrule
\end{tabular}
\end{table}

The required $G_t^{(\mathrm{req})}$ has several notable properties.

\begin{proposition}[Properties of $G_t^{(\mathrm{req})}$]\label{prop:Gt_props}
\begin{enumerate}
  \item[\rm(i)] $G_t^{(\mathrm{req})} > 0$ for $r < r_{\mathrm{cross}}$
    and $G_t^{(\mathrm{req})} < 0$ for $r > r_{\mathrm{cross}}$, where
    $r_{\mathrm{cross}} \approx 3.60\,\rs$ is the root of the cubic
    $r^3-4r^2\rs+3r\rs^2-3\rs^3 = 0$.
  \item[\rm(ii)] $G_t^{(\mathrm{req})} \to 0$ as $r \to \rs^+$ (horizon).
  \item[\rm(iii)] $G_t^{(\mathrm{req})} \to 0$ as $r \to \infty$ (flat space).
  \item[\rm(iv)] Near the horizon: $G_t^{(\mathrm{req})} \sim (r-\rs)/(r-\rs)^2
    \propto 1/(r-\rs)$, a single pole --- integrable.
\end{enumerate}
\end{proposition}

The sign change at $r_{\mathrm{cross}} \approx 3.6\,\rs$ is physical: $G_t$ is
repulsive near the horizon (pushing $\sigma_t$ outward) and attractive
in the asymptotic region (pulling $\sigma_t$ back toward zero).  This is
consistent with a self-stabilising field configuration.

\subsection{What this establishes and what remains open}

Theorem~\ref{thm:Gt} shows the three-term literature formula is incomplete.
The exact $G_t$ is known from~\eqref{eq:Gt_required}.  The open question is:
\emph{does a complete derivation of $G_\mu$ from the full action variation
reproduce~\eqref{eq:Gt_required} exactly?}  If yes, the Schwarzschild
$\sigma$-field is self-consistent and the moble criticism is fully answered.
If no, either the action must be extended (Path~2 of~\cite{matshaba2026:kin_var})
or the $G_\mu$ formula must be revised.

%=============================================================
\section{Result II: The TUO--QGD Architecture}
\label{sec:architecture}
%=============================================================

\subsection{The question}

TUO~\cite{matshaba2026:tuo} and QGD operate on overlapping domains.
TUO derives the Hot Big Bang initial conditions from two axioms (flat
Minkowski background + zero-sum charge constraint).  QGD derives spacetime
geometry from the Dirac field.  A natural question is whether TUO's axioms
are derivable from QGD, or whether the two theories are complementary.

\subsection{Layer 1: Axiom~I is a degenerate QGD solution}

\begin{proposition}[TUO Axiom~I from QGD]\label{prop:axiom1}
The QGD master metric with no sources, $\sigma_\mu = 0$, gives
$g_{\mu\nu} = \eta_{\mu\nu}$ (flat Minkowski).  TUO Axiom~I
(flat Minkowski background) is therefore the zero-source solution
of QGD.
\end{proposition}

This is immediate from the master metric
$g_{\mu\nu} = T^\alpha{}_\mu T^\beta{}_\nu
[\eta_{\alpha\beta} - \varepsilon\sigma_\alpha\sigma_\beta - \ldots]$:
setting $\sigma_\mu = 0$ recovers $\eta_{\mu\nu}$ exactly.
The important point is not that it is trivial but that it is
\emph{consistent}: TUO's starting geometry is not in conflict with QGD.

\subsection{Layer 2: $T_{\mathrm{TUO}}$ as a QGD horizon condition}

In QGD, the graviton field at a Schwarzschild horizon satisfies
$\sigt\big|_{r=\rs} = 1$ (this is the event horizon condition $\Sigma = 1$
in the master metric).  At the Planck scale with $M = M_\mathrm{Pl}$ and
$r = \lPl$:
\begin{equation}
  \sigt\bigg|_{\lPl,\,M_\mathrm{Pl}}
  = \sqrt{\frac{2GM_\mathrm{Pl}}{c^2\lPl}}
  = \sqrt{2} \approx 1.41.
\end{equation}
The Schwarzschild radius of a Planck mass is $\rs(M_\mathrm{Pl}) = 2\lPl$,
so the horizon condition $\sigt = 1$ is met at $r = 2\lPl$.

\begin{proposition}[TUO emergence temperature from QGD horizon condition]
\label{prop:TTUO}
The TUO result $\TTUO = (15/\pi^2)^{1/4}\TPl$ is the temperature of
a Planck-scale thermal bath whose Schwarzschild radius equals $2\lPl$.
This is the QGD condition $\sigt(r = 2\lPl,\, M_\mathrm{Pl}) = 1$,
i.e.\ the Planck-mass thermal bath sits exactly at its own QGD event horizon.
\end{proposition}

This is not a derivation of $\TTUO$ from QGD --- it connects the TUO
numerical result to a geometric condition that QGD makes precise.

\subsection{Layer 3: TUO as the pre-thermalisation transient}

\begin{table}[h]
\centering
\caption{Timeline connecting TUO and QGD.}
\label{tab:timeline}
\begin{tabular}{@{}p{2.8cm}p{4.5cm}p{4.5cm}@{}}
\toprule
Epoch & TUO description & QGD description \\
\midrule
$t < t_\mathrm{Pl}$ &
  Flat Minkowski background, no dynamics &
  $\sigma_\mu = 0$, $g_{\mu\nu} = \eta_{\mu\nu}$ \\[4pt]
$t \sim t_\mathrm{Pl}$ &
  Maximum fluctuation: all SM modes emerge simultaneously with zero-sum charges &
  Initial data $\sigma_t(t_\mathrm{Pl}) = \sqrt{2}$, $g_{\mu\nu} = \eta_{\mu\nu}$;
  Axiom~II imposed as constraint \\[4pt]
$t \in [t_\mathrm{Pl},\, 152\,t_\mathrm{Pl}]$ &
  Pre-thermalisation: flat-space expansion law $\ddot\sigma = c^2/\sigma$,
  QGP free-streaming &
  QGD master equation in flat/weakly curved background;
  $G_\mu, Q_\mu$ sourced by matter \\[4pt]
$t > 152\,t_\mathrm{Pl}$ &
  Junction to FRW (HBB); QGP thermalised &
  QGD attractor in FRW: $\dot\sigma_t = v$,
  $H \to H_\mathrm{FRW}$;
  $\rho_\sigma \approx \rho_\Lambda$ \\[4pt]
$t \sim H_0^{-1}$ &
  Not described by TUO (structure formation era) &
  QGD dark energy attractor;
  $\rho_\sigma = 3H_0^2c^4/8\pi G$ \\
\bottomrule
\end{tabular}
\end{table}

The QGD flat-space temporal equation (homogeneous $\sigma_t(t)$, no spatial
gradients) takes the form
\begin{equation}\label{eq:qgd_flat_temporal}
  \ddot\sigma_t = \sigma_t\,\dot\sigma_t^2.
\end{equation}
The TUO expansion law is $\ddot\sigma = c^2/\sigma$.  These are different
equations: the QGD flat-space equation is a genuinely nonlinear ODE with no
obvious analytic solution, while TUO's equation comes from energy-momentum
balance in Minkowski.  However, they describe the same physical epoch through
different variables: TUO uses the expansion radius $\sigma$ directly, while
QGD uses the dimensionless graviton phase.

\subsection{Layer 4: Axiom~II selects a unique QGD trajectory}

\begin{theorem}[TUO Axiom~II as a QGD initial-condition selector]
\label{thm:axiom2}
QGD admits a continuous family of solutions parametrised by the initial
data $(\sigma_\mu(t_0), \partial_t\sigma_\mu(t_0))$.  Among these,
TUO Axiom~II (zero-sum conserved charges: $B = L = Q = 0$ at $t_0 = t_\mathrm{Pl}$)
selects a single trajectory.  This constraint is not derivable from the
QGD gravitational field equations alone; it is a thermodynamic condition
on the matter sector.  However, it is fully compatible with QGD: the zero-sum
condition is preserved by the QGD evolution because the master field equation
conserves all Noether charges.
\end{theorem}

\begin{proof}[Sketch]
The QGD action is gauge-invariant under $U(1)_Y \times SU(2)_L \times SU(3)_c$
at low energies; correspondingly, the Noether currents $j^\mu_k$ are conserved.
The zero-sum condition $\int j^0_k\,d^3x = 0$ at $t_0$ is preserved under
evolution by $\partial_\mu j^\mu_k = 0$.  The constraint does not arise from
geometry; it is a statement about the state of the matter fields.
\end{proof}

\subsection{The unified picture}

\begin{figure}[h]
\centering
\renewcommand{\arraystretch}{1.4}
\begin{tabular}{@{}p{4cm}|p{4cm}|p{4cm}@{}}
  & \textbf{TUO role} & \textbf{QGD role} \\
\hline
Flat background & Axiom I & $\sigma_\mu = 0$ solution \\
Emergence temp.\ & $(15/\pi^2)^{1/4}\TPl$ & QGD horizon at $r = 2\lPl$ \\
Expansion law & $\ddot\sigma = c^2/\sigma$ & Flat-space transient \\
Zero-sum & Axiom II & IC constraint, not geometry \\
Thermalisation & $\tau_\mathrm{th} \approx 152\,t_\mathrm{Pl}$ & TUO--QGD junction \\
Dark energy & Not addressed & $\rho_\sigma \approx \rho_\Lambda$ at attractor \\
\end{tabular}
\caption{Unified TUO--QGD description.  TUO provides initial conditions;
QGD provides the evolution.}
\label{fig:unified}
\end{figure}

The key structural statement:

\begin{corollary}[TUO--QGD unification]
TUO and QGD are not competing theories.  They address different questions:
\begin{itemize}
  \item \emph{TUO asks:} What are the initial conditions of the universe,
    and why these rather than others?  Answer: the unique maximum fluctuation
    in a zero-sum Planck-scale thermal bath.
  \item \emph{QGD asks:} How does gravity work at all scales, from quantum
    to cosmological?  Answer: via the dynamics of the $\sigma$-field.
\end{itemize}
Together they form a complete description: TUO pins the initial data;
QGD evolves it.  The Hot Big Bang is the specific QGD solution selected
by the TUO initial conditions.
\end{corollary}

%=============================================================
\section{Discussion}
%=============================================================

\subsection{On the $G_t$ gap}

The three-term Kerr--Schild formula~\eqref{eq:Gt_formula} is not sufficient
to close the self-consistency gap.  There are two interpretations.

\emph{Interpretation A:} The $G_\mu$ formula in the QGD literature is
schematic; the full derivation from the action variation generates additional
terms that together produce~\eqref{eq:Gt_required}.  This is the most likely
scenario: the action variation tracks $\partial g^{\mu\nu}/\partial\sigma_t$,
$\partial\sqrt{-g}/\partial\sigma_t$, and the Christoffel dependence on
$\sigma_\mu$, all of which contribute terms that are not captured by the
simple Kerr--Schild directional derivative.

\emph{Interpretation B:} The action must be extended with a potential
$V(\sigma_t) = \sigma_t^2/8 - \sigma_t^4/16$, as identified
in~\cite{matshaba2026:kin_var}.  This potential emerges naturally from the
NJL effective theory and would close the gap with no free parameters.
Under this interpretation, the Schwarzschild $\sigma$-field is a
constrained solution, and the potential contributes to the $G_t$ balance.

Both interpretations are testable by a single explicit calculation: vary the
full QGD action with respect to $\sigma_t$ in the Schwarzschild background,
retain all $g(\sigma)$-dependence, and check whether the result
equals~\eqref{eq:Gt_required}.

\subsection{On the TUO--QGD relationship}

The most important structural result is that TUO Axiom~II (zero-sum charges)
is not a constraint on geometry but on the matter sector.  It cannot be
derived from QGD's gravitational dynamics alone.  This is not a weakness of
either theory; it reflects the fact that gravity and thermodynamics are
physically independent inputs.

The analogy to statistical mechanics is apt: the Boltzmann equation
(dynamics) does not determine the initial distribution function (initial
conditions).  TUO specifies the initial distribution; QGD evolves it.

The junction time $\tau_\mathrm{th} \approx 152\,t_\mathrm{Pl}$ from TUO
provides the handoff between the two theories: before it, TUO's flat-space
expansion law applies; after it, QGD's FRW master equation takes over.
Computing the QGD $\sigma$-field at $t = \tau_\mathrm{th}$ and verifying
it matches the TUO junction conditions would constitute a formal proof of
consistency between the two frameworks.

%=============================================================
\section{Summary of New Results}
%=============================================================

\begin{table}[h]
\centering
\caption{All new results.  Every numerical claim is reproduced by the
companion Python script (\texttt{qgd\_Gt\_tuo\_bridge.py}).}
\label{tab:summary}
\renewcommand{\arraystretch}{1.45}
\begin{tabular}{p{5.5cm}p{7cm}}
\toprule
Result & Statement \\
\midrule
\rowcolor{lightgray}
$G_t$ three-term formula is insufficient &
$G_t^{(\mathrm{lit})} \ne G_t^{(\mathrm{req})}$; residual $\Delta G_t$
given exactly by Eq.~\eqref{eq:Gt_gap} \\

Exact $G_t$ uniquely determined &
$G_t^{(\mathrm{req})} = -\sqrt{\rs}(r^3-4r^2\rs+3r\rs^2-3\rs^3)/
[4r^{7/2}(r-\rs)^2]$, with sign change at $r \approx 3.6\,\rs$ \\

\rowcolor{lightgray}
TUO Axiom~I from QGD &
$\sigma_\mu = 0$ gives $g_{\mu\nu} = \eta_{\mu\nu}$ exactly \\

$\TTUO$ as QGD horizon condition &
$(15/\pi^2)^{1/4}\TPl$ corresponds to $\sigt = 1$ at $r = 2\lPl$;
Planck-mass thermal bath at its own event horizon \\

\rowcolor{lightgray}
TUO Axiom~II as IC selector &
Zero-sum constraint selects unique QGD trajectory; not derivable from
geometry but preserved by QGD evolution \\

TUO = IC selector; QGD = dynamics &
Unified description: TUO fixes $(\sigma_\mu(t_\mathrm{Pl}),\, \Sigma Q_k=0)$;
QGD master equation evolves forward \\

\rowcolor{lightgray}
Junction time is the TUO--QGD handoff &
$\tau_\mathrm{th} \approx 152\,t_\mathrm{Pl}$ (TUO FRW cooling integral)
marks transition from TUO transient to QGD FRW phase \\
\bottomrule
\end{tabular}
\end{table}

%=============================================================
\section{Open Calculations}
%=============================================================

\begin{enumerate}
  \item \textbf{Full action variation for $G_\mu$:} Vary the complete QGD
    action in the Schwarzschild background, tracking all $g(\sigma)$
    dependence.  Show the result equals $G_t^{(\mathrm{req})}$.
    This closes the self-consistency proof.

  \item \textbf{TUO--QGD junction:} Compute the QGD $\sigma$-field at
    $t = \tau_\mathrm{th} = 152\,t_\mathrm{Pl}$ (the thermalisation time
    from TUO) and verify it matches the TUO junction conditions:
    $T = T_\mathrm{FRW}$, $H = H_\mathrm{FRW}$, $w = 1/3$.

  \item \textbf{QGD flat-space ODE analysis:} Find the general solution of
    $\ddot\sigma_t = \sigma_t\dot\sigma_t^2$ and characterise its
    relationship to the TUO expansion law $\ddot\sigma = c^2/\sigma$.

  \item \textbf{Zero-sum from SM symmetry:} Show that the SM gauge group
    $U(1)_Y \times SU(2)_L \times SU(3)_c$ requires $B-L = Q_\mathrm{em} = 0$
    at the Planck scale, deriving TUO Axiom~II from the Standard Model
    rather than postulating it.
\end{enumerate}

%=============================================================
\section{Conclusion}
%=============================================================

Two results have been established.

First, the exact $G_t$ required for self-consistency of the Schwarzschild
$\sigma$-field is uniquely determined: $G_t = -\sqrt{\rs}(r^3-4r^2\rs
+3r\rs^2-3\rs^3)/[4r^{7/2}(r-\rs)^2]$.  The three-term Kerr--Schild
formula in the QGD literature falls short; the residual gap is given
exactly.  The outstanding calculation---full action variation for $G_\mu$
in the Schwarzschild background---has a known target.

Second, TUO and QGD are complementary, not competing.  TUO is an
initial-condition theory; QGD is a dynamical theory.  TUO Axiom~I (flat
Minkowski) is the zero-source QGD solution.  TUO Axiom~II (zero-sum charges)
is a thermodynamic constraint that selects a unique trajectory in QGD
phase space.  The TUO emergence temperature $\TTUO = (15/\pi^2)^{1/4}\TPl$
corresponds to the QGD Schwarzschild-horizon condition at the Planck scale.
The thermalisation time $\tau_\mathrm{th} \approx 152\,t_\mathrm{Pl}$ marks
the handoff from TUO dynamics to QGD FRW evolution.

Together: TUO pins the initial data; QGD provides the law of motion.
The Hot Big Bang is the specific QGD solution selected by the TUO
constraints.

\begin{thebibliography}{9}

\bibitem{matshaba2026}
R.~Matshaba,
\textit{Universal Origins: Quantum Gravitational Dynamics},
University of South Africa (2026).
DOI: \href{https://doi.org/10.5281/zenodo.18827993}{10.5281/zenodo.18827993}.

\bibitem{matshaba2026:kin_var}
R.~Matshaba,
\textit{Kinetic Term Variation in QGD: Self-Consistency of the Schwarzschild
Solution}, companion paper, March 2026, same DOI.

\bibitem{matshaba2026:tuo}
R.~Matshaba,
\textit{TUO: Open Problems and Numerical Analysis},
companion paper, March 2026, same DOI.

\bibitem{matshaba2026:cosmo}
R.~Matshaba,
\textit{QGD Chapter~5: Cosmology (updated March 2026)},
companion paper, same DOI.

\bibitem{Damour1999}
T.~Damour and P.~Jaranowski,
Phys.\ Rev.\ D \textbf{61}, 064010 (2000).

\end{thebibliography}

\end{document}
